\documentclass[a4paper]{article}

%------------------------------------------------
% الحزم
%------------------------------------------------

\usepackage{enumitem}
\usepackage[margin=2.5cm]{geometry}
\usepackage{amsmath,amsfonts,amssymb}
\usepackage[charter]{mathdesign}
\usepackage{graphicx}
\usepackage{array}
\usepackage{xcolor}
\usepackage{needspace}
\usepackage{fancyhdr}
\usepackage{polyglossia}

%------------------------------------------------
% اللغات والخطوط
%------------------------------------------------

\setdefaultlanguage[
    calendar=gregorian,
    numerals=maghrib
]{arabic}

\setotherlanguage{english}

% إذا لم تكن الخطوط مثبتة في الجهاز، تُؤخذ من المجلد ../../fonts/
\IfFontExistsTF{Amiri}{%
    \newfontfamily\arabicfont[Script=Arabic, Scale=1.2]{Amiri}
}{%
    \newfontfamily\arabicfont[
        Script=Arabic, Scale=1.2,
        Path=../../fonts/, Extension=.ttf,
        UprightFont=*-Regular, BoldFont=*-Bold,
        ItalicFont=*-Italic, BoldItalicFont=*-BoldItalic
    ]{Amiri}
}

\IfFontExistsTF{AlBattar}{%
    \newfontfamily\arabicfontsf[Script=Arabic, Scale=1.3]{AlBattar}
    \newfontfamily\khatt[Script=Arabic, Scale=1.3]{AlBattar}
}{%
    \newfontfamily\arabicfontsf[
        Script=Arabic, Scale=1.3, Path=../../fonts/,
        BoldFont=ae_AlBattar.ttf
    ]{ae_AlBattar.ttf}
    \newfontfamily\khatt[
        Script=Arabic, Scale=1.3, Path=../../fonts/,
        BoldFont=ae_AlBattar.ttf
    ]{ae_AlBattar.ttf}
}

%------------------------------------------------
% نسخة الطالب / نسخة التصحيح
% لإنتاج نسخة التصحيح: عرّف \withanswers قبل \input لهذا الملف
%------------------------------------------------

\newif\ifcorrection
\ifdefined\withanswers \correctiontrue \else \correctionfalse \fi

\newcommand{\emptybox}{$\square$}
\ifcorrection
    \newcommand{\correctbox}{\textcolor{green!50!black}{$\boxtimes$}}
    \newcommand{\correct}[1]{\textcolor{green!40!black}{\textbf{#1}}}
    \newcommand{\explanation}[1]{%
        \par\vspace{3pt}
        \fcolorbox{gray!60}{gray!8}{\parbox{\dimexpr\linewidth-2\fboxsep-2\fboxrule}{%
        \small\textbf{التعليل:} #1}}}
\else
    \newcommand{\correctbox}{$\square$}
    \newcommand{\correct}[1]{#1}
    \newcommand{\explanation}[1]{}
\fi

\newcommand{\pts}[1]{\textbf{(#1 ن)}}

%------------------------------------------------
% إعداد الصفحة
%------------------------------------------------

\renewcommand{\footrulewidth}{1pt}
\renewcommand{\headrulewidth}{1pt}

\rhead{\khatt المدرسة العليا للأساتذة ببوزريعة}
\lhead{\khatt ملحقة بومرداس}

\cfoot{\thepage}
\lfoot{\khatt الأستاذ: د. بن عبد الرزاق يوسف}

\parindent 0pt
\setlength{\headsep}{15pt}
\arraycolsep=1.4pt

\pagestyle{fancy}

\setlist[enumerate,1]{label=\textbf{السؤال \arabic*:}, leftmargin=*, itemsep=12pt, align=left}

%================================================
% بداية الوثيقة
%================================================

\begin{document}

\begin{center}

{\Large\khatt\textbf{الواجب الأول}\ifcorrection\ \textbf{-- التصحيح النموذجي}\fi}

\vspace{0.2cm}

{\Large\khatt\textbf{في الذكاء الاصطناعي}}

\end{center}

\rule{\textwidth}{1pt}

\ifcorrection\else
\textbf{الاسم واللقب:} \dotfill \quad \textbf{الفوج:} \makebox[3cm]{\dotfill}

\vspace{0.3cm}
\fi

\begin{center}

{\khatt\textbf{أسئلة متعددة الاختيارات (20 سؤالا -- 15 نقطة)}}

\end{center}

\rule{\textwidth}{1pt}

\textbf{التعليمة:} ضع علامة \textbf{$\times$} في المربع المقابل للإجابة الصحيحة. لكل سؤال إجابة صحيحة واحدة فقط.

\vspace{0.3cm}

\begin{enumerate}

\needspace{9\baselineskip}\item
\textbf{أثناء كتابة رسالة، لاحظت أستاذة أن لوحة مفاتيح هاتفها تقترح الكلمة التالية باستمرار، ثم اقترح عليها تطبيق الخرائط أسرع طريق إلى المدرسة. قال لها زميلها إن الخاصيتين تعتمدان على المجال نفسه من مجالات الحوسبة. أي وصف يعبّر بدقة عن طبيعة هذا المجال وعن الهدف الذي يسعى إليه؟} \hfill \pts{0,75}
\begin{itemize}[leftmargin=1.6cm, labelsep=0.3cm, itemsep=2pt, topsep=4pt]
    \item[\emptybox\ \textbf{أ)}] هو مجموعة من قواعد البيانات الضخمة على الإنترنت تخزّن مسبقًا كل الجمل والطرق الممكنة، فيبحث الجهاز فقط عن إجابة جاهزة كتبها مهندس من قبل، دون أي تعلّم أو فهم أو تكيّف مع الوضعية الخاصة للشخص الذي يستعمل الهاتف أو التطبيق.
    \item[\emptybox\ \textbf{ب)}] هو تقنية عتادية مدمجة في المعالجات الحديثة تجعل الهواتف أسرع، وهذا ما يفسّر ظهور الاقتراحات فورًا؛ وجودة الاقتراح تتعلق أساسًا بمستوى البطارية ودقة الشاشة، لا ببرامج تحاكي قدرات الإنسان في الفهم أو في التفكير.
    \item[\correctbox\ \textbf{ج)}] \correct{هو فرع من علوم الحاسوب يجعل الحاسوب قادرًا على القيام بمهام تحتاج عادة إلى ذكاء الإنسان، مثل فهم اللغة والتعلّم من الأمثلة والإجابة عن الأسئلة، ولهذا تستطيع لوحة المفاتيح توقّع الكلمات ويستطيع تطبيق الخرائط اقتراح الطرق.}
    \item[\emptybox\ \textbf{د)}] هو فرع من اللسانيات يدرس كيف يتوقّع الناس الكلمات أثناء الكلام، ويكتفي صانعو الهواتف بنسخ القواعد التي اكتشفها اللسانيون في أجهزتهم، فالهاتف يطبّق قواعد نحوية ثابتة ولا يقوم بأي مهمة تشبه تفكير الإنسان أو تعلّمه.
\end{itemize}
\explanation{يُعرَّف هذا المجال بأنه فرع من علوم الحاسوب هدفه تمكين الآلة من أداء أنشطة تعتمد عادة على قدرات الإنسان الذهنية: الفهم والتعلّم والإجابة عن الأسئلة. والخدمات اليومية مثل اقتراح الكلمات وتخطيط المسارات أمثلة مألوفة على عمل هذا المجال، حتى وإن لم ينتبه إليها المستخدم.}

%------------------------------------------------

\needspace{9\baselineskip}\item
\textbf{سألت طالبة في قسم اللغة العربية: لماذا تُسمّى الأدوات المدروسة في القسم «توليدية» وليست «ذكية» فقط، مع أن محرّك البحث ذكي أيضًا؟ إذا انطلقنا من أصل كلمة «التوليدي» المأخوذة من الفعل «ولَّد» بمعنى أنشأ شيئًا، فما الذي يميّز هذه الفئة من الأدوات؟} \hfill \pts{0,75}
\begin{itemize}[leftmargin=1.6cm, labelsep=0.3cm, itemsep=2pt, topsep=4pt]
    \item[\correctbox\ \textbf{أ)}] \correct{هذه الأدوات تُنتج محتوى جديدًا، مثل فقرة أو صورة أو مقطع صوتي، لم يكن موجودًا بهذا الشكل من قبل، بدل الاكتفاء بالبحث عن وثائق مخزّنة مسبقًا وعرضها كما هي، وهو السلوك المعتاد لمحرّك البحث التقليدي.}
    \item[\emptybox\ \textbf{ب)}] هذه الأدوات تولّد حسابات موفّرة للطاقة، أي إنها تستهلك كهرباء أقل من محرّكات البحث؛ فالكلمة تشير إلى تقنيات توليد الطاقة المستعملة في مراكز البيانات التي تستضيفها، ولا علاقة لها بنوع المحتوى الذي يحصل عليه المستخدم.
    \item[\emptybox\ \textbf{ج)}] سُمّيت توليدية لأنها تنتمي إلى جيل جديد من الحواسيب ظهر بعد سنة ألفين، فالتسمية مصطلح تسويقي يصف عُمر الأجهزة التي تشغّلها، وليس وصفًا لطريقة استجابة البرنامج للطلب الذي يكتبه المستخدم.
    \item[\emptybox\ \textbf{د)}] هذه الأدوات تجمع أشهر صفحات الويب حول موضوع ما وتلصقها بعضها ببعض بالترتيب، فتكون النتيجة تجميعًا لنصوص موجودة؛ والكلمة تعني توليد قائمة من الروابط مرتّبة حسب عدد الزيارات التي حصلت عليها كل صفحة.
    \item[\emptybox\ \textbf{هـ)}] هذه الأدوات لا تُنتج إلا النصوص ولا تستطيع إنتاج الصور ولا الأصوات، وسُمّيت توليدية لأنها تولّد جملًا صحيحة نحويًا بتطبيق قواعد الصرف العربي التي برمجها يدويًا فريق من اللسانيين بالتعاون مع المطوّرين.
\end{itemize}
\explanation{السمة المميّزة لهذه العائلة من الأنظمة هي الإنشاء: فهي تنطلق مما يكتبه المستخدم لتصوغ مادة أصيلة في وسائط متعددة. أما نظام البحث فيحدّد مكان مادة مكتوبة سلفًا ويعيدها. وربط المصطلح بالفعل العربي الدال على إيجاد شيء جديد يساعد على تذكّر هذا الفرق.}

%------------------------------------------------

\needspace{9\baselineskip}\item
\textbf{رسم مكوّن على السبورة مخططًا بسيطًا من ثلاث خطوات ليوضّح كيف تعمل الأداة التوليدية منذ لحظة استعمالها، وطلب من أحد المشاركين ترتيب المراحل ترتيبًا صحيحًا، بدءًا من الشخص وانتهاءً بما يظهر على الشاشة. أي تسلسل يصف هذه العملية وصفًا صحيحًا؟} \hfill \pts{0,75}
\begin{itemize}[leftmargin=1.6cm, labelsep=0.3cm, itemsep=2pt, topsep=4pt]
    \item[\emptybox\ \textbf{أ)}] تُنتج الأداة أولًا محتوى عشوائيًا، ثم يقرؤه المستخدم ويكتب طلبًا يصف ما أعجبه منه، وأخيرًا تحلّل الأداة هذه الملاحظات لتقرّر هل تحتفظ بالمحتوى السابق أم تحذفه من سجلّ المحادثة.
    \item[\emptybox\ \textbf{ب)}] تبحث الأداة في الإنترنت عن صفحات موجودة، ثم يختار المستخدم إحدى الصفحات من قائمة النتائج، وبعدها تنسخ الأداة الصفحة المختارة كما هي، مع إضافة عنوان وبعض التنسيق قبل عرضها على المستخدم.
    \item[\emptybox\ \textbf{ج)}] يحمّل المستخدم مجموعة بيانات، ثم تعيد الأداة تدريب نفسها على هذه البيانات أثناء المحادثة، ولا تستطيع عرض إجابة توافق محتوى الملفات المحمّلة إلا بعد انتهاء إعادة التدريب التي قد تستغرق عدة ساعات.
    \item[\correctbox\ \textbf{د)}] \correct{يكتب المستخدم طلبًا يصف ما يحتاجه، ثم تحلّل الأداة هذا الطلب لتفهم المقصود منه وشروطه، ثم تُنتج محتوى جديدًا، كنص أو صورة، يستجيب لما طُلب في الطلب الأصلي.}
\end{itemize}
\explanation{تبدأ حلقة التفاعل الأساسية بمُدخل بشري، ثم تمرّ بمرحلة تأويل يعالج فيها النظام صياغة المُدخل ومقصده، وتنتهي بمُخرَج مصوغ حديثًا. وفهم هذا التسلسل يبيّن لماذا تهمّ جودة المُدخل الأول ووضوحه كثيرًا: فكل ما يليه يتوقّف على طريقة قراءة النظام له.}

%------------------------------------------------

\needspace{9\baselineskip}\item
\textbf{تريد أستاذة لغة عربية أن يستمع تلاميذها إلى قصيدة قديمة مقروءة مع ضبط الحركات ضبطًا صحيحًا، ليتدرّبوا على النطق في البيت. لا تملك الأستاذة أجهزة تسجيل ولا قارئًا محترفًا. بالنظر إلى أنواع المخرجات التي يستطيع الذكاء الاصطناعي التوليدي إنتاجها، أي نوع يلبّي حاجتها مباشرة؟} \hfill \pts{0,75}
\begin{itemize}[leftmargin=1.6cm, labelsep=0.3cm, itemsep=2pt, topsep=4pt]
    \item[\emptybox\ \textbf{أ)}] جدول تلخيصي يضع كل بيت من القصيدة في عمود ومعناه الحرفي في عمود آخر، لأن الجداول هي الشكل الرئيسي الذي تقدّمه الأدوات التوليدية للشعر، ويستطيع التلاميذ استنتاج النطق بقراءة المعاني بعناية في البيت.
    \item[\correctbox\ \textbf{ب)}] \correct{مُخرَج صوتي تُنتج فيه الأداة قراءة مسموعة للقصيدة مع مراعاة ضبط الحركات، لأن الأنظمة التوليدية قادرة على إنتاج الأصوات وليس النصوص المكتوبة فقط، مما يسمح بالاستماع إلى الأبيات دون الحاجة إلى قارئ بشري.}
    \item[\emptybox\ \textbf{ج)}] برنامج حاسوبي قصير يعدّ مقاطع كل بيت ويطبع اسم البحر الشعري، لأن البرامج هي المُخرَج الوحيد غير النصي الذي تستطيع هذه الأدوات إنتاجه، ومعرفة البحر تكفي التلاميذ لنطق الحركات نطقًا صحيحًا بمفردهم.
    \item[\emptybox\ \textbf{د)}] إعادة صياغة مكتوبة للقصيدة بالعربية الفصحى المعاصرة، لأن الأدوات التوليدية محصورة في النص ولا تتعامل مع الصوت، فعلى الأستاذة الاعتماد على نثر مبسّط يسهل على التلاميذ قراءته بصوت عالٍ دون نموذج مسموع.
    \item[\emptybox\ \textbf{هـ)}] رسم ثابت يصوّر المشهد الذي تصفه القصيدة، لأن الصور تساعد على الحفظ، والأدوات التوليدية تحوّل الشعر عادةً إلى صور لا إلى أصوات عند السؤال عن النطق، فيوجّه الرسم التلاميذ بطريقة غير مباشرة نحو الضبط الصحيح.
\end{itemize}
\explanation{لا تقتصر الأنظمة التوليدية على الكتابة؛ فمخرجاتها تشمل وسائط متعددة: نصوصًا وصورًا وأصواتًا وفيديو وجداول ومخططات وبرامج حاسوبية. ومن الضروري أن نختار نوع المُخرَج المناسب للهدف التعليمي: فإذا كان الهدف هو النطق، فإن مُخرَجًا يمكن الاستماع إليه يخدم المتعلّمين أفضل بكثير من أي بديل مكتوب.}

%------------------------------------------------

\needspace{9\baselineskip}\item
\textbf{خلال ورشة تكوينية، أُعطي المشاركون قائمة بمساعدات ذكاء اصطناعي توليدي مشهورة، وطُلب منهم ربط كل أداة بالشركة التي تطوّرها. تردّد أحد المشاركين أمام زوج واحد وطلب من المكوّن تأكيده. أي ربط من الروابط التالية بين الأداة والشركة المطوِّرة لها صحيح؟} \hfill \pts{0,75}
\begin{itemize}[leftmargin=1.6cm, labelsep=0.3cm, itemsep=2pt, topsep=4pt]
    \item[\correctbox\ \textbf{أ)}] \correct{تطوّر شركة \textenglish{Anthropic} أداة \textenglish{Claude}، وهي شركة متخصصة في الذكاء الاصطناعي تبني هذا المساعد وتحافظ عليه، وهو مثل الأدوات المشهورة الأخرى يقبل الطلبات المكتوبة بالعربية ويستطيع الإجابة بها للأساتذة والطلبة.}
    \item[\emptybox\ \textbf{ب)}] تطوّر شركة \textenglish{OpenAI} أداة \textenglish{Gemini}، وهي الشركة نفسها التي أنتجت \textenglish{ChatGPT}، وقد أطلقتها نسخةً خفيفة من مساعدها الرئيسي موجّهة أساسًا للهواتف وللمستخدمين الذين يفضّلون الإجابات القصيرة بعدة لغات منها العربية.
    \item[\emptybox\ \textbf{ج)}] تطوّر شركة \textenglish{Google} أداة \textenglish{Qwen} ضمن منتجات البحث، وقد صُمّمت خصيصًا للمستخدمين الصينيين والعرب الذين يريدون مساعدًا مرتبطًا مباشرة بصفحة نتائج البحث دون الحاجة إلى تطبيق مستقل أو حساب خاص.
    \item[\emptybox\ \textbf{د)}] تطوّر شركة \textenglish{Alibaba} أداة \textenglish{Perplexity}، وقد أنشأتها مساعدًا تجاريًا لأسئلة التسوّق عبر الإنترنت، وتُستعمل أساسًا لمقارنة أسعار المنتجات، لا لكتابة النصوص أو الإجابة عن الأسئلة التعليمية بالعربية أو بالإنجليزية.
\end{itemize}
\explanation{تتنافس عدة شركات في هذا المجال: \textenglish{OpenAI} تطوّر \textenglish{ChatGPT}، و\textenglish{Google} تطوّر \textenglish{Gemini}، و\textenglish{Anthropic} تطوّر \textenglish{Claude}، و\textenglish{DeepSeek} تطوّر مساعدًا يحمل اسمها، و\textenglish{Alibaba} تطوّر \textenglish{Qwen}، و\textenglish{Perplexity AI} تطوّر \textenglish{Perplexity}. ومعرفة الجهة التي تقف وراء كل أداة تساعد على إدراك أنها منتجات مختلفة لكل منها نقاط قوّته وأسلوبه.}

%------------------------------------------------

\needspace{9\baselineskip}\item
\textbf{أخبرت أستاذة زميلها أنها كتبت بضعة أسطر في مساعد ذكاء اصطناعي تطلب فيها أسئلة لفهم قصة قصيرة. فأجابها الزميل بأن لهذه الأسطر اسمًا خاصًا في مصطلحات هذا المجال، بالعربية وبالإنجليزية. ماذا يُسمّى النص الذي يكتبه المستخدم للأداة؟} \hfill \pts{0,75}
\begin{itemize}[leftmargin=1.6cm, labelsep=0.3cm, itemsep=2pt, topsep=4pt]
    \item[\emptybox\ \textbf{أ)}] يُسمّى المُخرَج (\textenglish{Output})، لأنه المعلومة التي تخرج من المستخدم وتتجه نحو الأداة؛ فالكلمة تشير إلى اتجاه الانتقال من الشخص إلى الآلة، لا إلى المحتوى الذي تعرضه الآلة لاحقًا على الشاشة.
    \item[\emptybox\ \textbf{ب)}] يُسمّى مجموعة البيانات (\textenglish{Dataset})، لأن كل رسالة يكتبها المستخدم تُضاف فورًا إلى مجموعة النصوص التي تعلّمت منها الأداة، والمصطلح يؤكّد أن المستخدم يساهم بمادة تدريب لا أنه يطلب مساعدة في مهمة معيّنة.
    \item[\emptybox\ \textbf{ج)}] يُسمّى الخوارزمية (\textenglish{Algorithm})، لأن المستخدم يعطي الحاسوب سلسلة من التعليمات البرمجية، ولا تقبل الأداة مُدخلًا صالحًا لتوليد أي إجابة إلا إذا كُتب بأسلوب رسمي متسلسل خطوة بخطوة.
    \item[\emptybox\ \textbf{د)}] يُسمّى الهلوسة (\textenglish{Hallucination})، لأن المستخدم يتخيّل الإجابة التي يريدها قبل كتابتها، فتحاول الأداة إعادة إنتاج تلك الإجابة المتخيَّلة؛ والمصطلح يصف الصورة الذهنية التي في ذهن المستخدم أثناء الكتابة.
    \item[\correctbox\ \textbf{هـ)}] \correct{يُسمّى الطلب أو الأمر، وبالإنجليزية \textenglish{Prompt}؛ وقد يكون سؤالًا أو تعليمة أو وصفًا للمهمة المطلوبة، وصياغته تؤثّر تأثيرًا كبيرًا في الإجابة التي تعيدها الأداة إلى المستخدم.}
\end{itemize}
\explanation{للرسالة التي يكتبها الشخص لتوجيه المساعد التوليدي اسم تقني خاص يقابل في العربية معنى الطلب أو التعليمة. ومعرفة هذا المصطلح مهمة، لأن جزءًا كبيرًا من مهارة استعمال هذه الأدوات يكمن في صياغة تلك الرسالة بعناية، باختيار كلمات توضّح المهمة والجمهور المستهدف والشكل المنتظر.}

%------------------------------------------------

\needspace{9\baselineskip}\item
\textbf{سأل طالب مساعد ذكاء اصطناعي عن قائل بيت شعري مشهور. أجاب المساعد بطلاقة وثقة، وذكر اسم شاعر وأضاف تاريخًا، لكن الطالب اكتشف لاحقًا في ديوان محقَّق أن البيت لشاعر آخر. ما المصطلح الذي يصف هذا السلوك لدى الأداة التوليدية؟} \hfill \pts{0,75}
\begin{itemize}[leftmargin=1.6cm, labelsep=0.3cm, itemsep=2pt, topsep=4pt]
    \item[\emptybox\ \textbf{أ)}] يُسمّى الإفراط في التعلّم، ويحدث عندما ترفض الأداة الإجابة عن أسئلة الشعر لأنها رأت قصائد كثيرة جدًا، فتذكر عمدًا شاعرًا آخر لتتجنّب تكرار الإجابة الأكثر شيوعًا في مادة تدريبها.
    \item[\correctbox\ \textbf{ب)}] \correct{يُسمّى الهلوسة، أي أن تُنتج الأداة معلومة خاطئة أو مختلَقة وتقدّمها بالثقة نفسها والأسلوب السلس نفسه اللذين تستعملهما مع المعلومات الصحيحة، فالنبرة الواثقة لا تضمن صحة المحتوى.}
    \item[\emptybox\ \textbf{ج)}] يُسمّى خطأ الترجمة، لأن الأداة تترجم دائمًا الأسئلة العربية إلى الإنجليزية داخليًا، فتُرجم اسم الشاعر خطأً عند إعادته إلى العربية، ويختفي الخطأ إذا طُرح السؤال مباشرة بالإنجليزية.
    \item[\emptybox\ \textbf{د)}] يُسمّى الرقابة، لأن الأداة مبرمجة لإخفاء أسماء بعض الشعراء القدامى لأسباب قانونية، فتضع اسمًا آخر مكانه، مع الإبقاء على تاريخ صحيح حتى تبدو الإجابة كاملة في نظر المستخدم.
\end{itemize}
\explanation{تصوغ الأنظمة التوليدية إجابات تبدو معقولة، والمعقولية ليست هي الحقيقة. فقد تنسب أبياتًا إلى غير قائليها، أو تختلق مراجع، أو تقدّم إعرابًا خاطئًا وهي تبدو واثقة تمامًا. وهذه الظاهرة هي السبب الرئيسي لوجوب التحقق من المعلومات المهمة في مصادر موثوقة قبل استعمالها في التدريس.}

%------------------------------------------------

\needspace{9\baselineskip}\item
\textbf{لكي يساعد أستاذ المبتدئين على فهم كيف تكتب الأداة التوليدية قصة أصيلة، شبّهها بتلميذ قرأ آلاف النصوص العربية فصار قادرًا بعد ذلك على كتابة نصوص جديدة بأسلوبه الخاص. ما الفكرة الأساسية التي يريد هذا التشبيه إيصالها عن هذه الأدوات؟} \hfill \pts{0,75}
\begin{itemize}[leftmargin=1.6cm, labelsep=0.3cm, itemsep=2pt, topsep=4pt]
    \item[\emptybox\ \textbf{أ)}] يبيّن التشبيه أن الأداة تحفظ كل نص حرفيًا، وعندما يُطلب منها قصة تسترجع قصة كاملة من ذاكرتها وتعيدها دون تغيير، تمامًا كما يسمّع التلميذ نصًا حفظه عن ظهر قلب أمام القسم.
    \item[\emptybox\ \textbf{ب)}] يبيّن التشبيه أن الأداة، مثل أي تلميذ، تحتاج إلى أستاذ حاضر في كل محادثة ليوافق على جملها، ولا تستطيع كتابة أي شيء ما لم يراجع إنسان كل كلمة قبل عرض الإجابة على المستخدم.
    \item[\emptybox\ \textbf{ج)}] يبيّن التشبيه أن الأداة لا تستطيع تقليد إلا أسلوب كاتب واحد في كل مرة، لأن التلميذ يُعجب عادة بكاتب واحد، فإذا طُلبت منها قصة بأسلوب مختلف ظهرت رسالة خطأ بدل نص جديد.
    \item[\correctbox\ \textbf{د)}] \correct{يبيّن التشبيه أن اطّلاع الأداة على كمية ضخمة من النصوص الموجودة يمكّنها من صياغة مادة جديدة تتبع أنماط اللغة المألوفة، كما تمكّن القراءةُ الواسعة التلميذَ من كتابة نصوص جديدة بدل نسخ ما قرأه.}
    \item[\emptybox\ \textbf{هـ)}] يبيّن التشبيه أن الأداة قرأت العدد نفسه من الكتب الذي يقرؤه تلميذ عادي، فمعرفتها محصورة في الكتب المدرسية، ولا تستطيع إنتاج محتوى عن مواضيع خارج المنهاج الرسمي الوطني.
\end{itemize}
\explanation{تجعل التشبيهات التقنية المجرّدة قريبة من الفهم. فصورة القارئ النهم الذي يصبح كاتبًا متمكّنًا توضّح أن هذه الأنظمة تبني قدرتها من كميات كبيرة جدًا من القراءة السابقة، ثم تصوغ جملًا أصيلة تتبع الأنماط التي استوعبتها. والفرق هو في الحجم: فالآلة عالجت نصوصًا أكثر بكثير مما يقرؤه أي إنسان.}

%------------------------------------------------

\needspace{9\baselineskip}\item
\textbf{ثُبّتت خاصيتان على لوحة إلكترونية في مدرسة: الأولى تقترح أقصر طريق بين قاعتين على خريطة المؤسسة، والثانية تكتب عند الطلب حوارًا قصيرًا جديدًا بين شخصيتين لنشاط لعب الأدوار. كيف نصنّف هاتين الخاصيتين بالنسبة إلى الذكاء الاصطناعي والذكاء الاصطناعي التوليدي؟} \hfill \pts{0,75}
\begin{itemize}[leftmargin=1.6cm, labelsep=0.3cm, itemsep=2pt, topsep=4pt]
    \item[\emptybox\ \textbf{أ)}] لا تُعدّ أي منهما من الذكاء الاصطناعي، لأن اقتراح الطريق مجرد حساب بسيط وكتابة الحوار مجرد تنسيق للنص؛ فالذكاء الاصطناعي لا يعني إلا الروبوتات الشبيهة بالإنسان التي تتحرك وتتفاعل مع الناس داخل القسم.
    \item[\emptybox\ \textbf{ب)}] خاصية الطريق وحدها من الذكاء الاصطناعي لأنها تقوم بحسابات، أما خاصية الحوار فهي معالجة نصوص عادية تشبه المدقق الإملائي، لأن كتابة النص لا تحتاج إلى أي قدرة من القدرات المرتبطة عادة بذكاء الإنسان.
    \item[\correctbox\ \textbf{ج)}] \correct{كلتاهما من الذكاء الاصطناعي لأن كل واحدة تؤدي مهمة تحتاج عادة إلى تفكير الإنسان، لكن كاتب الحوار وحده توليدي، لأنه يُنتج محتوى جديدًا استجابة لطلب بدل أن يقترح خيارًا موجودًا مسبقًا.}
    \item[\emptybox\ \textbf{د)}] كلتاهما توليدية، لأن أي خاصية تعرض شيئًا على الشاشة تُنتج محتوى جديدًا؛ فالطريق المرسوم على الخريطة صورة مولَّدة حديثًا، ولذلك لا يوجد فرق حقيقي بين الخاصيتين من حيث الصنف الذي تنتميان إليه.
\end{itemize}
\explanation{الأنظمة التوليدية جزء من مجال الذكاء الاصطناعي الأوسع. فكثير من التطبيقات الذكية تحلّل أو ترتّب أو توصي، وهذا وحده يحاكي التفكير البشري. أما ما يميّز الجزء التوليدي فهو أن مُخرَجه مادة مصوغة حديثًا، كنص أو صورة جديدين، مبنية على طلب المستخدم لا مختارة من خيارات موجودة.}

%------------------------------------------------

\needspace{9\baselineskip}\item
\textbf{أرسلت أستاذة الطلب نفسه تمامًا، وهو شرح المبني للمجهول مع أمثلة، إلى ثلاثة مساعدات ذكاء اصطناعي مختلفة، فجاءت الإجابات متباينة بوضوح في الطول والتنظيم والوضوح وعدد الأمثلة. ما الاستنتاج الأكثر حكمة والممارسة التي ينبغي اعتمادها في ضوء هذه الملاحظة؟} \hfill \pts{0,75}
\begin{itemize}[leftmargin=1.6cm, labelsep=0.3cm, itemsep=2pt, topsep=4pt]
    \item[\correctbox\ \textbf{أ)}] \correct{من الطبيعي أن تعطي الأدوات المختلفة إجابات مختلفة عن الطلب نفسه، لذلك من المفيد المقارنة بين مُخرَجات أكثر من مساعد، واختيار الإجابة، أو مزيج من عناصرها، الأنسب للمتعلّمين وللهدف التعليمي.}
    \item[\emptybox\ \textbf{ب)}] لا بد أن مساعدَين على الأقل معطّلان أو قديمان، لأن الأداة التي تعمل جيدًا تعطي دائمًا الإجابة نفسها التي تعطيها كل الأدوات الأخرى، فعلى الأستاذة الإبلاغ عن المشكلة واعتماد المساعد صاحب الإجابة الأطول فقط.
    \item[\emptybox\ \textbf{ج)}] يثبت هذا التباين أن الطلب كُتب باللغة الخطأ، لأن الطلبات المتطابقة المكتوبة بالعربية الصحيحة تعطي دائمًا إجابات متطابقة في كل الأدوات؛ فعلى الأستاذة إعادة كتابة الطلب بالإنجليزية قبل المحاولة من جديد.
    \item[\emptybox\ \textbf{د)}] على الأستاذة أن تختار دائمًا أقصر إجابة، لأن الإيجاز هو أوثق علامة على الدقة في الأدوات التوليدية، والإجابات الطويلة تحتوي عادة على هلوسات أكثر يصعب اكتشافها عند تحضير الدرس بسرعة.
    \item[\emptybox\ \textbf{هـ)}] الفروق شكلية فقط وتأتي من إعدادات الشاشة في كل موقع، أما المحتوى الفعلي للإجابات الثلاث فهو متطابق في العمق، لذا فالمقارنة بينها مضيعة للوقت لا تضيف شيئًا إلى تحضير الدرس.
\end{itemize}
\explanation{لكل مساعد تصميمه وتدريبه وأسلوبه، لذلك تؤدي التعليمة نفسها إلى إجابات تختلف في العمق والبنية والأمثلة. واعتبار هذا التنوع موردًا مفيدًا عادة مهنية جيدة: فوضع عدة مسوّدات جنبًا إلى جنب يكشف نقاط القوة والضعف، ويساعد المستخدم على انتقاء المادة الأنسب أو الجمع بين عناصرها.}

%------------------------------------------------

\needspace{9\baselineskip}\item
\textbf{كتبت أستاذة أولًا طلبًا قصيرًا تطلب فيه من الأداة شرح النعت. لم تكن راضية عن النتيجة، فأعادت كتابته لتطلب شرحًا موجّهًا لتلميذ في السنة الأولى متوسط، بلغة بسيطة، مع ثلاثة أمثلة من الحياة اليومية. لماذا تؤدي الصيغة الثانية عادة إلى إجابة أكثر فائدة؟} \hfill \pts{0,75}
\begin{itemize}[leftmargin=1.6cm, labelsep=0.3cm, itemsep=2pt, topsep=4pt]
    \item[\emptybox\ \textbf{أ)}] لأن الطلبات الأطول تُعالَج على خادم أقوى مخصّص للمستخدمين المدقّقين، فإضافة كلمات إلى أي طلب ترفع تلقائيًا جودة الإجابة، سواء أحملت الكلمات المضافة معلومات مفيدة عن المهمة أم لم تحمل.
    \item[\emptybox\ \textbf{ب)}] لأن الأداة تتجاهل الطلب الأول في أي محادثة ولا تجيب إلا ابتداءً من الرسالة الثانية، فالتحسّن يأتي من ترتيب الرسائل، لا من أي تغيير في مضمون ما كُتب في الطلب.
    \item[\emptybox\ \textbf{ج)}] لأن ذكر المستوى الدراسي يفعّل نمطًا خاصًا ينسخ الكتاب المدرسي الرسمي حرفيًا، مما يضمن الصحة، أما التفاصيل المتعلقة باللغة والأمثلة فتتجاهلها الأداة فعليًا عند صياغة إجابتها.
    \item[\correctbox\ \textbf{د)}] \correct{لأنها تحدّد المستوى المستهدف ونوع اللغة وعدد الأمثلة، فتعطي الأداة قيودًا واضحة حول الجمهور والشكل، فتأتي الإجابة مفصّلة على مقاس ما تحتاجه الأستاذة فعلًا بدل أن تكون شرحًا عامًا.}
\end{itemize}
\explanation{التعليمة الغامضة تترك للنظام أن يخمّن الجمهور ومستوى اللغة ومقدار التفصيل، فيميل إلى إجابة عامة. أما إضافة عناصر محددة، مثل مستوى المتعلّم وبساطة اللغة وعدد دقيق من الأمثلة، فتضيّق الاحتمالات وتوجّه المُخرَج نحو الحاجة الحقيقية في القسم. فوضوح المُدخَل ودقته يحسّنان النتيجة مباشرة.}

%------------------------------------------------

\needspace{9\baselineskip}\item
\textbf{تلقّت أستاذة متربّصة فقرة مولَّدة بالذكاء الاصطناعي تتضمّن قاعدة نحوية تنوي تقديمها في القسم. ذكّرها مؤطّرها بخطوات التحقق التي دُرست في المقياس. أي ترتيب للخطوات يعكس هذه الطريقة في التحقق من معلومة مهمة قبل استعمالها؟} \hfill \pts{0,75}
\begin{itemize}[leftmargin=1.6cm, labelsep=0.3cm, itemsep=2pt, topsep=4pt]
    \item[\emptybox\ \textbf{أ)}] تسأل الأداة نفسها هل إجابتها صحيحة، وتقبل الإجابة إذا أكّدتها الأداة، ثم تعيد صياغة القاعدة بأسلوبها، وأخيرًا تقدّمها للتلاميذ، لأن تأكيدًا ثانيًا من الأداة يساوي الرجوع إلى مصدر خارجي.
    \item[\emptybox\ \textbf{ب)}] تقرّر فورًا هل أعجبتها الإجابة، ثم تبحث عن مصدر يوافق قرارها، وتقارن الأجزاء المتطابقة فقط، وأخيرًا تقرأ بقية الإجابة إذا سمح الوقت بذلك قبل بداية الحصة.
    \item[\correctbox\ \textbf{ج)}] \correct{تقرأ الإجابة، وتحدّد المعلومة التي يجب التحقق منها، وتبحث عنها في مصدر موثوق، وتقارن ما يقوله المصدر بما قالته الأداة، ثم تقرّر هل تستعمل المعلومة أو تصحّحها أو ترفضها.}
    \item[\emptybox\ \textbf{د)}] تنسخ الإجابة في مساعد ثانٍ ثم في ثالث، وتحتفظ بالصيغة الأكثر تكرارًا، لأن اتفاق عدة أدوات هو الخطوة الأخيرة ويُغني عن الرجوع إلى الكتب والمعاجم أو أي مرجع بشري آخر.
    \item[\emptybox\ \textbf{هـ)}] تبحث أولًا عن مصدر موثوق دون قراءة الإجابة، ثم تقرأ الإجابة بسرعة، وتقرّر بحسب درجة الثقة التي كتبت بها الأداة، لأن النبرة الواثقة أوضح دليل على أن المعلومة مأخوذة من ذلك المصدر الموثوق.
\end{itemize}
\explanation{يتبع التحقق المسؤول تسلسلًا مقصودًا: قراءة متأنية، ثم عزل المعلومة المحددة، ثم الرجوع إلى مرجع يُعتمد عليه، ثم المقابلة بين الصيغتين، وبعدها فقط إصدار الحكم. وتخطّي بعض الخطوات أو الاعتماد على الأداة لتتحقق من نفسها يُفرغ العملية من معناها، لأن النظام الذي أنتج الخطأ قد يكرّره بالثقة نفسها.}

%------------------------------------------------

\needspace{9\baselineskip}\item
\textbf{تحضّر أستاذة لغة عربية درس الأسبوع القادم حول نص قراءة، وتريد توفير الوقت دون التضحية بالجودة. تفكّر في عدة طرق لإشراك مساعد ذكاء اصطناعي توليدي في تحضيرها. أي استعمال من الاستعمالات التالية يقدّمه المقياس طريقةً مناسبةً لاستفادة أستاذ اللغة من هذه الأداة؟} \hfill \pts{0,75}
\begin{itemize}[leftmargin=1.6cm, labelsep=0.3cm, itemsep=2pt, topsep=4pt]
    \item[\correctbox\ \textbf{أ)}] \correct{أن تطلب من الأداة صياغة مسوّدة أسئلة لفهم نص القراءة، ثم تقرؤها وتصحّح ما يلزم وتكيّفها مع مستوى تلاميذها قبل استعمالها في القسم.}
    \item[\emptybox\ \textbf{ب)}] أن تترك الأداة تصحّح كل تعابير التلاميذ الكتابية وترسل العلامات مباشرة إلى الإدارة دون مراجعتها، لأن التصحيح الآلي أكثر موضوعية من أي أستاذ ولا يحتاج إلى أي تدقيق إضافي.
    \item[\emptybox\ \textbf{ج)}] أن تطلب من الأداة تقديم الدرس كاملًا للتلاميذ وحدها عبر جهاز العرض بينما تغادر هي القسم، لأن الأداة تعرف القواعد النحوية أفضل، والتلاميذ ينتبهون للآلة أكثر مما ينتبهون للإنسان.
    \item[\emptybox\ \textbf{د)}] أن تستعمل الأداة بديلًا كاملًا عن الكتاب المدرسي الرسمي، فتوزّع إجاباتها على أنها المرجع المعتمد في النحو والأدب، لأن المحتوى المولَّد أحدث وبالتالي أدق من أي مادة مدرسية مطبوعة.
    \item[\emptybox\ \textbf{هـ)}] أن تنقل أول إجابة تقدّمها الأداة عن النص مباشرة إلى السبورة، لأن إعادة قراءة المادة المولَّدة أو تعديلها يُذهب الوقت الذي يبرّر استعمال المساعد أصلًا.
\end{itemize}
\explanation{هذه المساعدات مفيدة لإنتاج مسوّدات: أمثلة جديدة لقاعدة نحوية، وأسئلة لفهم النصوص، وشروح مبسّطة لدروس البلاغة، وموضوعات للتعبير الكتابي. لكن إسهامها نقطة انطلاق فقط؛ فالمربّي يبقى يراجع ويصحّح ويختار ما يناسب متعلّميه، فلا يكون كسب الوقت على حساب الحكم المهني أو الدقة.}

%------------------------------------------------

\needspace{9\baselineskip}\item
\textbf{تحتاج أستاذة إلى التأكد من الإعراب الصحيح لكلمة ومن القائل الحقيقي لبيت شعري قبل إدراجهما في اختبار، فوضعت قائمة بالمراجع التي يمكن الرجوع إليها. أي مما يلي لا يُعدّ مصدرًا موثوقًا للتحقق من هذا النوع من المعلومات؟} \hfill \pts{0,75}
\begin{itemize}[leftmargin=1.6cm, labelsep=0.3cm, itemsep=2pt, topsep=4pt]
    \item[\emptybox\ \textbf{أ)}] معجم عربي معروف صادر عن هيئة أكاديمية معترف بها، يُعدّ مداخلَه متخصصون، ويستعمله الأساتذة والباحثون على نطاق واسع للتأكد من معاني الكلمات واستعمالها الصحيح.
    \item[\correctbox\ \textbf{ب)}] \correct{إجابة ثانية يولّدها المساعد نفسه بعد إعادة السؤال عليه، لأن ردًّا متكرّرًا من الأداة التي ربما أنتجت الخطأ في المرة الأولى لا يصلح تأكيدًا مستقلًا لما ادّعته.}
    \item[\emptybox\ \textbf{ج)}] كتاب نحو معتمد يُدرَّس في الجامعات، ألّفه علماء متمكّنون في النحو العربي وروجع عبر طبعات عديدة، يشرح القواعد ويقدّم أمثلة معربة يعتمد عليها الأساتذة كثيرًا.
    \item[\emptybox\ \textbf{د)}] ديوان شاعر محقَّق تحقيقًا علميًا، أعدّه باحث قابل بين المخطوطات ووثّق نسبة كل قصيدة إلى صاحبها، فهو مرجع يُعتمد عليه للتأكد من قائل بيت معيّن.
    \item[\emptybox\ \textbf{هـ)}] الكتاب المدرسي الرسمي المعتمد من وزارة التربية، الذي راجعه مفتّشو المادة، وهو المرجع المشترك الذي يستعمله الأساتذة والتلاميذ طوال السنة الدراسية.
\end{itemize}
\explanation{لا ينجح التحقق إلا إذا كان المرجع مستقلًا عن المادة التي نتحقق منها، وراجعه أشخاص مؤهَّلون. والمعاجم المعروفة، وكتب النحو المعتمدة، والدواوين المحقَّقة، والكتب المدرسية الرسمية تستوفي هذا الشرط. أما إعادة سؤال المولِّد نفسه فلا تستوفيه، لأنه قد يكرّر خطأه بالثقة ذاتها.}

%------------------------------------------------

\needspace{9\baselineskip}\item
\textbf{يحضّر رئيس قسم عرضًا لزملائه، ويريد مخططًا يقارن نتائج التلاميذ عبر ثلاثة فصول، إضافة إلى برنامج قصير يحسب المعدلات تلقائيًا من جدول بيانات. شكّ زميل في قدرة المساعد التوليدي على المساعدة في أي من المهمتين. كيف ينبغي أن يردّ رئيس القسم؟} \hfill \pts{0,75}
\begin{itemize}[leftmargin=1.6cm, labelsep=0.3cm, itemsep=2pt, topsep=4pt]
    \item[\emptybox\ \textbf{أ)}] الزميل محقّ، لأن الأدوات التوليدية محصورة في إنتاج فقرات نثرية، وكل ما يتعلق بالأرقام أو المخططات أو البرمجة يجب إنجازه يدويًا أو ببرامج إحصاء متخصصة لا علاقة لها بالذكاء الاصطناعي.
    \item[\emptybox\ \textbf{ب)}] تستطيع الأداة إنتاج المخطط دون البرنامج، لأن إنشاء البرامج الحاسوبية حكر على المطوّرين المحترفين، والمساعدات التوليدية ترفض أي طلب يحتوي كلمات تقنية مثل صيغة أو برنامج أو جدول بيانات.
    \item[\emptybox\ \textbf{ج)}] تستطيع الأداة كتابة البرنامج دون المخطط، لأن المخرجات المرئية تتطلب كاميرا لالتقاط صور حقيقية، والأنظمة التوليدية لا تستطيع إنشاء أي عنصر مرئي لم يصوّره شخص من قبل.
    \item[\correctbox\ \textbf{د)}] \correct{تستطيع المساعدات التوليدية المساعدة في المهمتين، لأن مخرجاتها تشمل الجداول والمخططات والبرامج الحاسوبية، إضافة إلى النصوص والصور والأصوات والفيديو، مع ضرورة مراجعة النتائج قبل عرضها على الزملاء.}
\end{itemize}
\explanation{نطاق المواد التي تستطيع هذه الأنظمة إنشاءها واسع: نصوص مكتوبة، وصور، وكلام وأصوات، وفيديو، وجداول ومخططات منظَّمة، وبرامج حاسوبية. ومعرفة هذا الاتساع تسمح للمربّين باستعمال مساعد واحد في مهام تحضير كثيرة، مع تذكّر أن المخرجات الرقمية والتقنية تستحق المراجعة الدقيقة نفسها التي تستحقها أي إجابة مكتوبة.}

%------------------------------------------------

\needspace{9\baselineskip}\item
\textbf{عرض تلميذ على أستاذه إعرابًا لجملة ولّده الذكاء الاصطناعي. كان الإعراب منسّقًا تنسيقًا جميلًا، ويستعمل مصطلحات نحوية دقيقة، ومكتوبًا بثقة تامة، فأصرّ التلميذ على أنه صحيح لا محالة. أي استدلال ينبغي أن يعتمده الأستاذ في الرد على حجة التلميذ؟} \hfill \pts{0,75}
\begin{itemize}[leftmargin=1.6cm, labelsep=0.3cm, itemsep=2pt, topsep=4pt]
    \item[\emptybox\ \textbf{أ)}] التلميذ محقّ، لأن الأدوات التوليدية لا تستعمل المصطلحات المتخصصة إلا عندما تكون متأكدة، فالمفردات الدقيقة علامة مدمجة تدلّ على أن الإعراب قد تمّ التحقق منه داخليًا في كتب النحو المعتمدة.
    \item[\emptybox\ \textbf{ب)}] التلميذ محقّ فقط إذا كانت الإجابة أطول من نصف صفحة، لأن الإعراب الطويل يدل على أن الأداة درست احتمالات أكثر، والإجابات القصيرة هي الأكثر احتمالًا لاحتواء معلومات مختلَقة أو خاطئة.
    \item[\emptybox\ \textbf{ج)}] التلميذ مخطئ فقط لأن الأداة استُعملت بالعربية؛ فالإعراب نفسه لو طُلب بالإنجليزية لكان صحيحًا حتمًا، لأن الأخطاء النحوية لا تظهر إلا عندما تعمل الأداة بلغات غير الإنجليزية.
    \item[\emptybox\ \textbf{د)}] التلميذ مخطئ لأن الإجابة المكتوبة بإتقان خاطئة دائمًا، فالأدوات تضيف تنسيقًا أنيقًا خصيصًا لإخفاء الأخطاء، ولذلك يجب رفض كل إجابة تبدو مصقولة دون أي تحقق إضافي.
    \item[\correctbox\ \textbf{هـ)}] \correct{العرض الجذّاب والواثق لا يثبت الصحة، لأن هذه الأدوات قد تقدّم إعرابًا خاطئًا بالثقة نفسها التي تقدّم بها الصحيح؛ فيجب التحقق من الإعراب في مرجع نحوي موثوق قبل قبوله أو استعماله.}
\end{itemize}
\explanation{الأسلوب والمضمون مستقلان في المحتوى المولَّد. فقد يُنتج النظام نصًا سلسًا منظّمًا بمصطلحات لا عيب فيها، ومع ذلك يرتكب أخطاء معرفية أو نحوية. ويجب تدريب المتعلّمين على الفصل بين مدى إقناع الإجابة ومدى صحتها، وعلى اعتبار المراجع الموثوقة الحَكَم النهائي في الصحة.}

%------------------------------------------------

\needspace{9\baselineskip}\item
\textbf{في اجتماع للأساتذة، قال أحد الزملاء إن المدارس لن تحتاج قريبًا إلى أساتذة اللغة العربية، ما دامت المساعدات التوليدية تستطيع إعداد الأمثلة والأسئلة والشروح في ثوانٍ. وفق الرؤية التي يقدّمها المقياس، ما الموقف الأدق من دور الأستاذ؟} \hfill \pts{0,75}
\begin{itemize}[leftmargin=1.6cm, labelsep=0.3cm, itemsep=2pt, topsep=4pt]
    \item[\emptybox\ \textbf{أ)}] سيبقى الأستاذ ضروريًا فقط لضبط الانضباط داخل القسم، أما كل القرارات البيداغوجية واختيار المحتوى والتصحيح فسيتولّاها المساعد، لأنه أوثق من الحكم البشري في مسائل اللغة.
    \item[\correctbox\ \textbf{ب)}] \correct{الأداة مساعد وليست بديلًا: يبقى الأستاذ هو المسؤول عن القراءة والتصحيح واختيار ما يناسب تلاميذه، لأن المحتوى المولَّد قد يحتوي أخطاء ولا يعرف خصوصيات القسم واحتياجاته.}
    \item[\emptybox\ \textbf{ج)}] على الأساتذة تجنّب الأدوات التوليدية تمامًا، لأن أي استعمال لها ينقل المسؤولية تدريجيًا إلى الآلة، والطريقة الوحيدة للحفاظ على دورهم المهني هي تحضير كل المواد يدويًا دون أي مساعدة رقمية.
    \item[\emptybox\ \textbf{د)}] الزميل محقّ في النحو والإملاء اللذين تتقنهما الآلات إتقانًا تامًا، لكن الأساتذة سيبقون ضروريين للشعر، لأن الأدوات التوليدية لا تستطيع معالجة الشعر ولا إنتاج أي محتوى يتعلق بالأدب القديم.
\end{itemize}
\explanation{تسرّع هذه الأنظمة التحضير لكنها لا تستطيع تحمّل مسؤولية ما يُدرَّس. فالحكم المهني، ومعرفة المتعلّمين، وواجب التحقق من الدقة، كلها تبقى عند المربّي. واعتبار التقنية أداة مساعدة يُبقي الخبير البشري في المركز، فهو الذي يقرّر ما هو صحيح ومناسب ومفيد لفوج معيّن من التلاميذ.}

%------------------------------------------------

\needspace{9\baselineskip}\item
\textbf{تريد أستاذة تبسيط درس في البلاغة حول التشبيه لتلاميذ السنة الثانية ثانوي. طلبها الحالي يقول فقط: اشرح التشبيه. بالنظر إلى ما يجعل الطلب فعّالًا، أي صيغة معدَّلة هي الأقرب إلى إنتاج إجابة تطابق حاجتها داخل القسم؟} \hfill \pts{0,75}
\begin{itemize}[leftmargin=1.6cm, labelsep=0.3cm, itemsep=2pt, topsep=4pt]
    \item[\emptybox\ \textbf{أ)}] اشرح التشبيه بتفصيل كبير وبأكبر قدر ممكن من الشمول، مع كل الآراء التاريخية للبلاغيين القدامى، وكل الأقسام، والنقاش العلمي كاملًا، حتى لا يفوت شيء ويختار التلاميذ ما يريدون قراءته.
    \item[\emptybox\ \textbf{ب)}] أرجوك، أرجوك اشرح التشبيه جيدًا جدًا وصحيحًا جدًا، لأن هذا الدرس مهم للغاية لتلاميذي، وأنا أحتاج حقًا إلى أفضل إجابة ممكنة اليوم دون أي خطأ على الإطلاق.
    \item[\correctbox\ \textbf{ج)}] \correct{اشرح التشبيه لتلاميذ السنة الثانية ثانوي بلغة عربية واضحة وبسيطة، وعرّف أركانه الأربعة، وقدّم مثالين من الشعر المشهور ومثالين من الكلام اليومي، في جدول قصير يليه ملخّص موجز.}
    \item[\emptybox\ \textbf{د)}] شرح التشبيه الآن، قصير، سريع، أمثلة، جدول ربما، بسيط، تلاميذ، ثانوي، شكرًا؛ وبهذا يبقى الطلب موجزًا فتخمّن الأداة بقية التفاصيل بنفسها وتقرّر المستوى والشكل الأنسب.
    \item[\emptybox\ \textbf{هـ)}] اشرح التشبيه كما ورد في الكتاب المدرسي، لكن لا تذكر مستوى التلاميذ ولا عدد الأمثلة، لأن إضافة مثل هذه المعلومات تقيّد الأداة وتقلّل من إبداع شرحها.
\end{itemize}
\explanation{التعليمات الفعّالة تحدّد الجمهور، ومستوى اللغة، والمضمون المنتظر بدقة، وعدد الأمثلة، والشكل المطلوب. أما التأدّب أو الإطالة أو سرد كلمات مفتاحية مبتورة فلا يوفّر هذا التوجيه. وعندما تكون هذه العناصر صريحة، لا يحتاج النظام إلى التخمين، وتنسجم إجابته أكثر بكثير مع الهدف التعليمي.}

%------------------------------------------------

\needspace{9\baselineskip}\item
\textbf{لاحظت تلميذة أنها عندما طلبت من مساعد تلخيص قصة قصيرة، ثم أعادت الطلب بصياغة مختلفة قليلًا حدّدت فيها جمهورًا أصغر سنًا، تغيّر التلخيص في مفرداته وطوله. فتساءلت هل الأداة غير موثوقة. كيف يُفسَّر هذا السلوك على أفضل وجه؟} \hfill \pts{0,75}
\begin{itemize}[leftmargin=1.6cm, labelsep=0.3cm, itemsep=2pt, topsep=4pt]
    \item[\correctbox\ \textbf{أ)}] \correct{الأداة تستجيب لصياغة كل طلب، فتغيير العبارة أو إضافة تفاصيل مثل الجمهور المستهدف يغيّر الإجابة بشكل طبيعي؛ وهذه الحساسية متوقَّعة ويمكن استغلالها عمدًا للحصول على نتيجة أنسب.}
    \item[\emptybox\ \textbf{ب)}] يدل التغيير على أن الأداة نسيت القصة بين الطلبين واختلقت قصة جديدة، وهذا يعني أن كل إجابة تأتي بعد الطلب الأول في المحادثة يجب اعتبارها هلوسة والتخلّص منها فورًا.
    \item[\emptybox\ \textbf{ج)}] الأداة تُنتج عمدًا إجابات مختلفة لتختبر انتباه المستخدم، والإجابة الأولى هي الصحيحة دائمًا، أما الصيغ اللاحقة فتُضعَّف عمدًا لتثبيط المستخدمين عن طرح السؤال نفسه مرتين.
    \item[\emptybox\ \textbf{د)}] سبب الاختلاف هو سرعة الاتصال بالإنترنت فقط، لأن الاتصالات البطيئة تتلقّى تلخيصات أقصر؛ أما صياغة الطلب فلا تأثير لها في المحتوى ولا في الطول ولا في المفردات التي تُنتجها الأداة.
\end{itemize}
\explanation{يتشكّل المُخرَج المولَّد مباشرة بحسب صياغة المُدخَل. فذكر جمهور مختلف أو طول أو نبرة مختلفة يجعل النظام يعدّل المفردات والعمق والبنية. وهذه الاستجابة ليست عيبًا، بل هي الرافعة التي يملكها المستخدم لتوجيه النتائج: فإعادة الصياغة المدروسة من أنجع الطرق للحصول على إجابة أفضل.}

%------------------------------------------------

\needspace{9\baselineskip}\item
\textbf{تخطّط أستاذة مبتدئة لاستعمال الذكاء الاصطناعي التوليدي في إعداد ورقة عمل كاملة حول أنواع الخبر لتلاميذ المتوسط. طلبت من مؤطّرتها طريقة عمل شاملة تجمع أهم دروس المقياس في صياغة الطلب، والأدوات، والدقة، والمسؤولية. أي خطة تعكس هذه التوجيهات أفضل من غيرها؟} \hfill \pts{0,75}
\begin{itemize}[leftmargin=1.6cm, labelsep=0.3cm, itemsep=2pt, topsep=4pt]
    \item[\emptybox\ \textbf{أ)}] تكتب طلبًا من كلمة واحدة في أداة واحدة، وتطبع الإجابة الأولى فورًا لتوفير الوقت، ولا تصحّح الأخطاء إلا إذا اشتكى التلاميذ لاحقًا، لأن التحقق المسبق من المحتوى يُلغي الفعالية التي دفعت إلى استعمال الأداة.
    \item[\emptybox\ \textbf{ب)}] تطرح السؤال الغامض نفسه على ثلاث أدوات، وتحتفظ بالإجابة الأطول لأنها حتمًا الأكمل، وتوزّعها دون تعديل، معتمدة على أن اتفاق الأدوات في النقاط العامة يضمن صحة كل التفاصيل.
    \item[\emptybox\ \textbf{ج)}] تتجنّب كتابة أي تفاصيل في الطلب حتى تبقى الأداة حرّة في إبداعها، ثم تسأل الأداة نفسها هل تحتوي ورقة العمل على أخطاء، وتعتمد تأكيدها خطوةَ تحقق نهائية قبل توزيع الورقة على التلاميذ.
    \item[\correctbox\ \textbf{د)}] \correct{تكتب طلبًا دقيقًا يحدّد المستوى واللغة وعدد الأمثلة، وتقارن مخرجات أكثر من أداة، وتتحقق من المحتوى النحوي في مرجع موثوق، ثم تكيّف ورقة العمل وتضع لمستها النهائية بنفسها قبل توزيعها على التلاميذ.}
    \item[\emptybox\ \textbf{هـ)}] تترك التلاميذ يستعملون الأداة مباشرة لإنشاء أوراق عملهم في البيت، لأن ذلك ينقل مسؤولية الدقة إليهم، ولا يحتاج الأستاذ إلا إلى جمع الأوراق دون مراجعة صحة القواعد والأمثلة.
\end{itemize}
\explanation{تربط طريقة العمل السليمة بين العادات الأساسية في هذا الفصل: صياغة تعليمة واضحة ومحددة، والاستفادة من تنوّع المساعدات المتاحة، وتأكيد المعلومات المهمة في مراجع موثوقة، وإبقاء المربّي مسؤولًا عن المنتَج النهائي. فكل خطوة تعالج حدًّا معروفًا لهذه الأدوات، ومجموعها يحوّل مسوّدة سريعة إلى مادة تعليمية موثوقة.}

\end{enumerate}

\vfill
\begin{center}
{\khatt بالتوفيق}
\end{center}

\end{document}
